\chapter{Análise por modelo}

\eyebrow{segmentação de gordura no lote}

\vspace{4pt}
A visão \textbf{Analysis} de um lote executa os modelos treinados sobre suas imagens e
transforma os resultados em um espaço de trabalho científico: uma galeria, uma tabela ordenável com uma
distribuição, conformação, correlação física e comparação lado a lado.

\section{Executando a análise}

Abra um lote e selecione \textbf{Analysis}. Pressione \tele{Analyze batch} para processar
todas as imagens ainda não analisadas, com uma barra de progresso ao vivo. Os resultados são salvos no
lote, de modo que reabri-lo não os executa novamente; \tele{Reanalyze} força uma nova execução.
O botão \tele{CSV} exporta todos os resultados.

Cada imagem passa pelo mesmo pipeline:

\begin{enumerate}[leftmargin=1.4em,itemsep=3pt]
  \item \textbf{Remoção de fundo.} A carcaça é recortada e composta sobre
        preto --- as mesmas condições em que o modelo de gordura foi treinado.
  \item \textbf{Segmentação de gordura.} O modelo validado produz uma máscara de gordura; a
        fração de gordura é medida \emph{em relação à superfície da carcaça}, não
        à foto inteira.
  \item \textbf{Conformação.} A medida de convexidade roda sobre a silhueta limpa
        (Estação~8).
  \item \textbf{Grau (opcional).} Com \tele{include grade (experimental)}
        marcado, o grau de acabamento experimental também é calculado.
\end{enumerate}

\begin{caution}[title={Por que a remoção de fundo importa}]
O modelo de gordura foi treinado com carcaças recortadas sobre fundo preto. Sem a
remoção, ele conta ladrilhos claros e o jaleco do operador como gordura (dando valores
falsamente altos). A remoção é automática; uma pílula \tele{background removed} a confirma,
e o percentual de gordura passa a ser a fração da \emph{superfície da carcaça}.
\end{caution}

\section{Galeria}

A galeria (Figura~\ref{fig:analysis-gallery}) mostra cada carcaça com sua sobreposição
de gordura em ciano e o percentual de gordura abaixo --- uma varredura visual rápida do
lote inteiro. Acima dela, estatísticas do lote (média, desvio padrão, mediana, mín, máx) e uma
linha de pílulas de status: \pillok{validated segmentation $\cdot$ IoU 0.92},
\pillalert{experimental grade $\cdot$ n=22} e \tele{analysis saved to batch}.

\shot{analysis-gallery.jpeg}{Analysis $\rightarrow$ Gallery. Sobreposição de gordura (ciano) por
carcaça com a fração de gordura; estatísticas do lote e as pílulas de status codificadas por cor
acima; o Inspector informa a média, o desvio padrão e a amplitude de gordura do lote.}{analysis-gallery}

\section{Tabela e distribuição}

A aba \textbf{Table + distribution} (Figura~\ref{fig:analysis-table}) lista as
carcaças em uma tabela ordenável (por etiqueta ou \% de gordura), ao lado do grau e das
colunas de referência física, junto a um histograma da distribuição de gordura do lote.

\shot{analysis-table.jpeg}{Analysis $\rightarrow$ Table + distribution. Uma tabela ordenável
de gordura, grau e medidas físicas por carcaça, ao lado do histograma de gordura do
lote.}{analysis-table}

\section{Comparar}

A aba \textbf{Compare} (Figura~\ref{fig:analysis-compare}) permite selecionar até
quatro carcaças e ver suas sobreposições e números lado a lado --- útil para
calibrar a percepção contra o modelo.

\shot{analysis-compare.jpeg}{Analysis $\rightarrow$ Compare. Até quatro carcaças
lado a lado, cada uma com sua sobreposição de gordura, percentual de gordura e grau experimental.}{analysis-compare}

\section{Correlação física}

A aba \textbf{Physical correlation} plota a \% de gordura do modelo contra a referência
física que você inseriu (espessura de gordura, área de olho de lombo), com o $r$ de Pearson. Esta
é a aba que vai \emph{fechar o ciclo de validação} assim que o conjunto de dados pareado for
grande o suficiente: ela testa se a medida de superfície prediz a referência
física medida no animal.
